\begin{table*}[t]
\centering
\footnotesize
\begin{tabular}{l S[table-format=1.0] S[table-format=3.0] c S[table-format=1.3] S[table-format=1.3] S[table-format=2.2] r}
\toprule
Unit & {$k$} & {Units} & Trunc. & {Mean} & {Chance} & {Ratio} & Perm. $p$ \\
\midrule
items & 3 & 425 & 0/90 & 0.144 & 0.014 & 10.23 & $<$ 0.00050 \\
 & 4 & 425 & 18/90 & 0.161 & 0.019 & 9.01 & $<$ 0.00050 \\
 & 5 & 425 & 36/90 & 0.187 & 0.024 & 8.91 & $<$ 0.00050 \\
\midrule
tasks & 3 & 52 & 18/90 & 0.315 & 0.115 & 2.87 & $<$ 0.00050 \\
 & 4 & 52 & 27/90 & 0.331 & 0.154 & 2.42 & $<$ 0.00050 \\
 & 5 & 52 & 42/90 & 0.369 & 0.192 & 2.28 & $<$ 0.00050 \\
\bottomrule
\end{tabular}
\caption{Transfer at every $k$ under both ranking units, over the same 425 clean items grouped into 52 tasks and the same ten frontier configurations, 90 ordered pairs each. A task's score is the share of its items that flip. Trunc. counts pairs where one side had fewer than $k$ non-zero units, so its top-$k$ was cut short rather than padded. Chance is recomputed per pair from the effective $k$ after truncation, which is why the coarser unit carries a chance line an order of magnitude higher and a correspondingly smaller ratio at the same permutation $p$. $p$ is over 2{,}000 permutations that shuffle each configuration's values among its own positions.}
\label{tab:tasklevel}
\end{table*}
